\input{../tex-inputs/latex-leseansicht-vorspann}

\section[Arthur Schnitzler an Gustav Schwarzkopf, 19. 7. 1898]{L04469 Arthur Schnitzler an Gustav Schwarzkopf, 19. 7. 1898}
\nopagebreak\mylabel{L04469v}
\rehead{ }\normalsize\beginnumbering\briefempfaengerindex{Schwarzkopf, Gustav@\textsc{Schwarzkopf, Gustav}!zzzSchnitzler, Arthur@\emph{von Arthur Schnitzler}!1898-07-192@{19. 7. 1898}|(be}
\toendnotes[C]{\smallbreak\pagebreak[2]}
\correspDesc{Versand  durch Arthur Schnitzler am 19. 7. 1898 in Bad Fusch
\newline{}Übermittlung  am 20. 7. 1898 in Bad Fusch
\newline{}Erhalt  durch Gustav Schwarzkopf im Zeitraum [20. 7. 1898 – 24. 7. 1898?] in Wien}\toendnotes[C]{\smallbreak}
\Standort{DLA, A:Schnitzler, HS.1985.1.1897.}
\physDesc{Bildpostkarte, 109 Zeichen
\newline{}Handschrift: Bleistift, lateinische Kurrent
\newline{}Versand: 1) Stempel: »\nobreak{}\oindex{Bad Fusch@\textbf{Bad Fusch}|pwk}Bad Fusch, 20/7 98\nobreak{}«.   2) Stempel: »\nobreak{}\oindex{I., Innere Stadt@\textbf{I., Innere Stadt}|pwk}Wien 1/1, 21. 7. 98, bestellt\nobreak{}«. }\toendnotes[C]{\smallbreak}\pstart{}{\pb}Herrn \textsc{Gustav Schwarzkopf}\pend{}\pstart{}Wien\oindex{Wien@\textbf{Wien}|pw}\pend{}\pstart{}\textsc{I. Tiefer Graben 23}\oindex{Wien@\textbf{Wien}!I., Innere Stadt@\textbf{I., Innere Stadt}!Tiefer Graben 23@\textbf{Tiefer Graben 23}|pw}\pend{}{\bigskip}
\pstart
           \noindent{}{\pb}\textcolor{gray}{\textbf{Gruss aus Bad Fusch\oindex{Bad Fusch@\textbf{Bad Fusch}|pw}.}}\pend
           \vspace{1em}
\pstart
           \noindent{}{\pb}von
               Dr. Hofmannstal\pwindex{Hofmannsthal, Hugo August von 21.\,12.\,1841 Wien – 8.\,12.\,1915 ebd.@\textsc{Hofmannsthal, Hugo August von} (21.\,12.\,1841 Wien – 8.\,12.\,1915 ebd.)|pw}, Frau Anna\pwindex{Hofmannsthal, Anna von 27.\,1.\,1849 Wien – 22.\,3.\,1904 Sanatorium Fürth@\textsc{Hofmannsthal, Anna von} (27.\,1.\,1849 Wien – 22.\,3.\,1904 Sanatorium Fürth)|pw} und \spacefill\mbox{Arthur Schn}\pend
           
\pstart
           \label{T_L04469-1v}\edtext{Dienſtag Abend {\\}19/7 98}{\lemma{\textnormal{\emph{Dienstag Abend 19/7 98}}}\Cendnote{\textnormal{Die Datierung ist entlang des rechten Randes geschrieben.}}}\label{T_L04469-1}\pend
           \selectlanguage{ngerman}\endnumbering\briefempfaengerindex{Schwarzkopf, Gustav@\textsc{Schwarzkopf, Gustav}!zzzSchnitzler, Arthur@\emph{von Arthur Schnitzler}!1898-07-192@{19. 7. 1898}|)be}\mylabel{L04469h}
\begin{anhang}
\end{anhang}\newcommand{\dateiname}{L04469}\newcommand{\titel}{Arthur Schnitzler an Gustav Schwarzkopf, 19. 7. 1898}\newcommand{\editorInnen}{Herausgegeben von Jahnke, SelmaMüller, Martin Anton}\input{../tex-inputs/latex-leseansicht-abspann}
